\documentclass[11pt,a4paper]{article}

% ---------------------------------------------------------------- page
\usepackage[a4paper,top=18mm,bottom=17mm,left=22mm,right=22mm,
            headsep=7mm,footskip=11mm]{geometry}
\usepackage[T1]{fontenc}
\usepackage{lmodern}
\usepackage{microtype}          % kerning + margin protrusion
\usepackage{enumitem}
\usepackage{tabularx}
\usepackage{fancyhdr}
\usepackage{titlesec}
\usepackage{xcolor}
\usepackage{fontawesome5}
\usepackage{academicons}
\usepackage[hidelinks]{hyperref}

\definecolor{accent}{HTML}{16365C}   % one accent, used only for links
\hypersetup{colorlinks=true, urlcolor=accent, linkcolor=accent,
            pdftitle={Duo Yang - Curriculum Vitae}, pdfauthor={Duo Yang},
            pdfsubject={Curriculum Vitae},
            pdfkeywords={natural language processing, tabular data, knowledge graphs,
                         table representation learning, retrieval-augmented generation}}

\setlength{\parindent}{0pt}
\setlength{\emergencystretch}{2em}
\raggedright                      % never justify a CV
\widowpenalty=10000 \clubpenalty=10000
\hyphenpenalty=10000 \exhyphenpenalty=10000  % ragged right: don't hyphenate

% ---------------------------------------------------------------- headings
% more space above than below: the header binds to what follows
\titleformat{\section}{\normalsize\bfseries\scshape}{}{0pt}{}
\titlespacing*{\section}{0pt}{12pt plus 3pt minus 3pt}{4pt}
\titleformat{\subsection}{\small\bfseries}{}{0pt}{}
\titlespacing*{\subsection}{0pt}{9pt plus 2pt minus 1pt}{3pt}

% ---------------------------------------------------------------- helpers
\newcommand{\cvdot}{\,\textperiodcentered\,}
\newcommand{\ic}[1]{\texorpdfstring{#1\,}{}}
\newcommand{\cvann}[1]{{\small #1}}
\newcommand{\cvlinks}[1]{{\small #1}}

% date rail on the left, content indented -- one alignment system throughout
\newlength{\raillen}\setlength{\raillen}{30mm}
\newcommand{\cvrow}[2]{%
  \noindent\parbox[t]{\raillen}{\small #1}%
  \parbox[t]{\dimexpr\linewidth-\raillen\relax}{#2}\par\vspace{3pt}}
\newcommand{\cvrowsub}[3]{%
  \noindent\parbox[t]{\raillen}{\small #1}%
  \parbox[t]{\dimexpr\linewidth-\raillen\relax}{#2\\[1pt]{\small #3}}\par\vspace{5pt}}

\pagestyle{fancy}\fancyhf{}
\renewcommand{\headrulewidth}{0pt}
\fancyhead[L]{\scriptsize\itshape Duo Yang}
\fancyhead[R]{\scriptsize\itshape Curriculum Vitae}
\fancypagestyle{first}{\fancyhf{}%
}

% ---------------------------------------------------------------- bibliography
\usepackage[backend=bibtex,style=numeric,sorting=ydnt,maxbibnames=99,
            giveninits=true,doi=true,url=false,defernumbers=true]{biblatex}
\AtEveryBibitem{\clearfield{issn}\clearfield{isbn}\clearfield{month}}
\addbibresource{cv.bib}

% bold the CV owner in every author list
% Bold my own name in every author list. \mkbibnamefamily never fires with
% the bibtex backend, and \ifstrequal does not expand \namepartfamily --
% \ifdefstring compares the macro's definition, which is what works here.
\DeclareNameFormat{default}{%
  \ifdefstring{\namepartfamily}{Yang}%
    {\textbf{\usebibmacro{name:given-family}
       {\namepartfamily}{\namepartgiveni}{\namepartprefix}{\namepartsuffix}}}%
    {\usebibmacro{name:given-family}
       {\namepartfamily}{\namepartgiveni}{\namepartprefix}{\namepartsuffix}}%
  \usebibmacro{name:andothers}}

\DeclareFieldFormat{titlecase}{#1}
\DeclareFieldFormat[inproceedings,inbook,article,incollection]{title}{\textbf{#1}}
\DeclareFieldFormat{booktitle}{\emph{#1}}   % venue in italics -- the one legit italic
\DeclareFieldFormat{journaltitle}{\emph{#1}}
\DeclareFieldFormat{doi}{\href{https://doi.org/#1}{\aiDoi}}
\renewcommand*{\bibfont}{\small}
\setlength{\bibitemsep}{3.5pt}
\setlength{\bibhang}{1.6em}                 % hanging indent

\begin{document}
\thispagestyle{first}

% ---------------------------------------------------------------- header
{\LARGE\bfseries\href{https://duo0301.github.io}{\textcolor{black}{Duo Yang}}}\quad
{\small\href{mailto:duo.yang@kuleuven.be}{duo.yang@kuleuven.be}}\hfill
{\small
\href{https://duo0301.github.io}{\faGlobe} \cvdot{}
\href{https://orcid.org/0009-0008-5942-3397}{\aiOrcid} \cvdot{}
\href{https://scholar.google.com/citations?user=apSK8IcAAAAJ}{\aiGoogleScholar} \cvdot{}
\href{https://github.com/duo0301}{\faGithub} \cvdot{}
\href{https://www.linkedin.com/in/duo-yang-614526231}{\faLinkedin}}\\[3pt]
{\normalsize PhD Researcher in Computer Science \cvdot{} KU Leuven}
\vspace{4pt}\hrule height 0.5pt
\vspace{2pt}

\section*{Research interests}
Natural language processing, tabular data and knowledge graphs, working towards
artificial intelligence for data integration. Particular interest in table
representation learning and the problems it raises, such as data scarcity, and in
retrieval-augmented generation for industrial use cases such as troubleshooting.

\section*{Experience}
\cvrowsub{2023 --\,now}{\textbf{Doctoral Researcher}\cvdot{}\href{https://dtai.cs.kuleuven.be/}{DTAI},
  \href{https://www.kuleuven.be/kuleuven}{KU Leuven}\cvdot{}E2E core lab,
  \href{https://www.kuleuven.be/flandersmake}{Flanders Make@KU Leuven}}
  {Table representation learning and semantic table annotation; knowledge graph
   construction and ontology engineering. Technical solutions for industrial partners,
   including retrieval-augmented generation systems for troubleshooting and
   ontology-driven root cause analysis. Advisor:
   \href{https://www.kuleuven.be/wieiswie/en/person/00142590}{Prof.\ Anastasia Dimou}.}

\section*{Education}
\cvrow{2023 --\,now}{\textbf{PhD in Computer Science}\cvdot{}\href{https://www.kuleuven.be/kuleuven}{KU Leuven}, Belgium}
\cvrow{2020--2022}{\textbf{MSc in IT \& Cognition}\cvdot{}\href{https://www.ku.dk/}{University of Copenhagen}, Denmark}
\cvrow{2014--2018}{\textbf{BEng in Telecommunication Engineering}\cvdot{}\href{https://en.szu.edu.cn/}{Shenzhen University}, China}

\section*{Publications}
{\small \textdagger\ equal contribution.\quad \textbf{Bold} marks my name.\quad
\faFilePdf\,paper \quad \faFilePowerpoint\,slides \quad \faImage\,poster \quad \faGlobe\,project site \quad \faDatabase\,data \quad \aiDoi\,DOI}
\vspace{2pt}
\nocite{*}
\printbibliography[keyword=conference, heading=subbibliography, title={Peer-reviewed conference papers}]
\printbibliography[keyword=workshop,   heading=subbibliography, title={Workshop and challenge papers}]
\printbibliography[keyword=demo,       heading=subbibliography, title={Demonstrations and posters}]

\section*{Awards and honours}
\cvrow{2024}{Best Poster Award, E2E cluster\cvdot{}Flanders Make Scientific Conference}
\cvrow{2023}{Accuracy track, \textbf{1st Prize}\cvdot{}SemTab 2023 challenge @ ISWC \href{https://sem-tab-challenge.github.io/2023/}{\faLink}}

\section*{Software and datasets}
\cvrowsub{}{\href{https://dtai-kg.github.io/MLSO/}{\textbf{MLSea / MLSO ontology}}}
  {Semantic layer and ontology for discoverable machine learning metadata.}
\cvrowsub{}{\href{https://rexpek.github.io/ishikawa-diagram-ontology/ishikawa-diagram-ontology/latest/index.html}{\textbf{Ishikawa Diagram Ontology}}}
  {Domain ontology supporting root cause analysis.}

\section*{Talks}
\cvrowsub{2026}{Label-Constrained Column Annotation with Language Models and Graph Neural Networks}
  {ICDE 2026, conference talk \cvdot{} \href{https://docs.google.com/presentation/d/1y3FvEmdQ3s1inufP0be86mznur9UYWVu/preview}{\faFilePowerpoint}}
\cvrowsub{2023}{TorchicTab: Semantic Table Annotation with Wikidata and Language Models}
  {SemTab challenge @ ISWC 2023, challenge talk \cvdot{} \href{https://docs.google.com/presentation/d/1A7MwTxjogUsiJzF5Qi6_Me5PJjj2fNy2/preview}{\faFilePowerpoint}}

\section*{Teaching}
\input{teaching.tex}
\CourseRows

\section*{Master's thesis supervision}
\ThesisRows

\section*{Availability}
Seeking a \textbf{research internship} in winter/spring 2026--2027, and
\textbf{full-time researcher or scientist positions} from July 2027.

\end{document}
